% !TEX root = ../../tree_notes.tex
% =====================================================================
%  Chapter 4 review -- the 15-section revision contract (chapter_end.txt).
%  Revision-only: tables, one-liners, pointers back via \thmref/\probref/
%  \secref/\coderef/\propref. No derivations, no re-teaching.
%  Summit: LC Hard 834 Sum of Distances in Tree (rooting + reroot shift).
% =====================================================================

\chaptersection{Chapter 4 Summary}

Three sections, one physical question: \emph{where do the edges live, and which node is
on top}. Storage (4.1) decides what a "child" is and how far navigation costs; rooting
(4.2) turns symmetric edges into a directed structure in one DFS; encodings (4.3) decide
what string the judge handed you and how to read it back. Every entry below is a pointer,
not a re-derivation.

\begin{center}
\renewcommand{\arraystretch}{1.25}
\begin{tabular}{@{}p{0.25\linewidth} p{0.47\linewidth} p{0.20\linewidth}@{}}
\toprule
\textbf{Result} & \textbf{Statement (with conditions)} & \textbf{Where} \\
\midrule
Four storage models & Pointer / adjacency / parent array / implicit index all implement
  \code{children(v)}; they differ in up-navigation and edge cost. & \defref{c04-forms} \\
Implicit index is exact iff complete & $2i{+}1,2i{+}2$ wastes no slot only for a complete
  tree; a skewed tree needs up to $2^{\height}$ slots. & \thmref{c04-implicit} \\
Parent array $\to$ children & One $O(n)$ loop turns \code{par[]} into \code{kids[]} lists.
  & \coderef{c04-reps} \\
Rooting is one DFS & A parent-guarded walk from $r$ visits each node once and yields the
  (order-independent) $\parent$, $\depth$, $\size$ arrays. & \thmref{c04-dfs-orient} \\
Size counts yourself & $\size[u]\ge 1$; a leaf has $\size=1$, and $\size[r]=n$.
  & \secref{c04-rooting} \\
Reroot across an edge & $f(c)=f(p)+(n-2\,\size[c])$ -- moving the root is a local,
  $O(1)$ edit. & \propref{c04-reroot} \\
Encoding must mark nulls & A value-only order is decodable uniquely only if nulls (or a
  second order) are written down. & \thmref{c04-selfdelim} \\
Level order is BFS-packed & The judge's array with \code{null}s is a breadth-first
  packing; read it with a \emph{queue}, not a preorder recursion. & \defref{c04-levelorder} \\
Depth-tagged preorder decodes & Dashes (or \# marks) give the depth, and a path stack of
  size $=$ depth rebuilds pointers in one pass. & \probref{c04-recover-preorder} \\
\bottomrule
\end{tabular}
\end{center}

\medskip
{\small\textbf{Algorithms from this chapter -- the \emph{input format} picks the reader.}}
\begin{center}
\small
\renewcommand{\arraystretch}{1.2}
\setlength{\tabcolsep}{3pt}
\begin{tabular}{@{}p{0.19\linewidth} l l p{0.29\linewidth} p{0.155\linewidth}@{}}
\toprule
\textbf{Algorithm} & \textbf{Time} & \textbf{Memory} & \textbf{When to use (deciding phrase)} & \textbf{Code} \\
\midrule
Children fetch, four models
  & $O(n)$ & $O(n)$
  & pick storage by which navigation you need
  & \coderef{c04-reps} \\
Root + read \code{par/dep/sz/ht}
  & $O(n)$ & $O(n){+}O(h)$
  & undirected edges, question needs above/below
  & \coderef{c04-root} \\
Centroid leaf-peel
  & $O(n)$ & $O(n)$
  & the root that minimises height
  & \coderef{c04-min-height} \\
Reroot two-pass
  & $O(n)$ & $O(n)$
  & the same aggregate for \emph{every} root
  & \coderef{c04-summit} \\
Level-order decode
  & $O(n)$ & $O(n)$
  & input is LeetCode's array with nulls
  & \coderef{c04-formats} \\
Preorder-with-depth decode
  & $O(L)$ & $O(n)$
  & input is a dashed / \#-marked preorder string
  & \coderef{c04-recover-preorder} \\
\bottomrule
\end{tabular}
\end{center}

\chaptersection{Master Decision Flowchart}

Post-study ladder: walk top to bottom, stop at the first branch that matches. It
\emph{decides} the representation and the reader -- a follower never opens the wrong file.

\begin{lstlisting}[language={}, basicstyle=\ttfamily\small, numbers=none, frame=single]
  How does the tree ARRIVE, and what must I navigate?
  |
  |-- it is already a pointer TreeNode / N-ary Node
  |     +-- just visit it -> that is a traversal, Chapter 3 [c03]
  |     +-- "go to parent" often -> rebuild a parent[] or pass p down [4.1]
  |
  |-- it is parallel arrays leftChild[]/rightChild[], or parent[]
  |     +-- need children down  -> one loop builds kids[]        [4.1, LC 1376]
  |     +-- "is this ONE valid tree?" -> in-degree<=1, one root,
  |         DFS reaches all n                                        [4.1, LC 1361]
  |
  |-- it is n-1 UNDIRECTED edges (LeetCode "rooted at 0", CF 1-indexed)
  |     +-- build adj both ways, then ROOT it: parent-guarded DFS  [4.2]
  |     +-- "count over subtree(v)" -> one bottom-up return        [4.2]
  |     +-- "best root / min height" -> centroid, peel leaves      [4.2, LC 310]
  |     +-- "answer for EVERY root" -> root once, then O(1) reroot [4.2, Summit 834]
  |
  |-- it is a STRING or a null-filled ARRAY
  |     +-- array with nulls -> level-order: read with a QUEUE     [4.3]
  |     +-- preorder with # or dashes -> read with a STACK         [4.3, LC 1028]
  |     +-- brackets  0(1(4)(5))(2)  -> recursive split; emit () for
  |         a right-only child                                       [4.3, LC 606]
  |
  +-- tree may be a CHAIN (n ~ 2e5) -> recursion is O(h) deep:
      root/decode iteratively (explicit stack), or sizes in reverse BFS order.
      Complete-tree heap layout -> Chapter 7 [c07]; subtree-as-range tin/tout -> [c24];
      answer from every root in full generality -> Rerooting [c20]; Prufer -> [c32].
\end{lstlisting}

\chaptersection{Common Mistakes Checklist}

\begin{itemize}[nosep, leftmargin=1.4em]
  \item \textbf{No parent guard on an undirected \code{adj}.} The DFS bounces straight
        back where it came -- infinite recursion (\secref{c04-reps}, Trap~1;
        \secref{c04-rooting}, Trap~1).
  \item \textbf{Assuming $2i{+}1$ indexing is compact for \emph{any} tree.} Only a
        complete tree packs without holes (\secref{c04-reps}, Trap~2; \thmref{c04-implicit}).
  \item \textbf{Conflating undirected $\degop$ with rooted child count.} A rooted node
        has $\degop-1$ children; the root has $\degop$ (\secref{c04-reps}, Trap~3).
  \item \textbf{A guard on the wrong array.} The skip test belongs on the traversal's
        parent parameter, not a stale index (\secref{c04-rooting}, Trap~1).
  \item \textbf{Subtree size that forgets the node itself.} Initialise $\size[u]=1$, not
        $0$; every size is otherwise one too small (\secref{c04-rooting}, Trap~2).
  \item \textbf{Recursing on a chain-shaped tree.} Depth $\approx n$ blows the stack and
        \emph{looks} like WA -- root iteratively (\secref{c04-rooting}, Trap~3).
  \item \textbf{Building a node for a \# / null token.} A marker is an \emph{absent}
        child; return \code{nullptr}, do not allocate (\secref{c04-formats}, Trap~1).
  \item \textbf{Reading a level-order array as preorder.} It is BFS-packed -- a queue,
        not a preorder recursion (\secref{c04-formats}, Trap~2).
  \item \textbf{Dropping the empty \code{()} for a right-only child.} Without it a lone
        child is re-read as a left child -- the ambiguity LC 606 keeps on purpose
        (\secref{c04-formats}, Trap~3; \thmref{c04-selfdelim}).
\end{itemize}

\chaptersection{WA Clinic}

Verdicts are symptoms; the causes are the short list this chapter drilled. Diagnose
top-down -- every leaf names the training problem.

\begin{figure}[H]
\centering
\begin{tikzpicture}[diagram,
    v/.style={boxnode, fill=primaryblue!12, draw=primaryblue, text width=22mm, align=center},
    b/.style={boxnode, fill=difforange!12, draw=difforange, text width=12mm},
    l/.style={boxnode, fill=successgreen!10, draw=successgreen, text width=58mm,
              font=\scriptsize, align=left}]
  \node[v] (verdict) {\textbf{Verdict on the judge}};
  \node[b, below left=8mm and 6mm of verdict] (wa) {\textbf{WA}};
  \node[b, below right=1mm of verdict] (tle) {\textbf{TLE}};
  \node[b, below left=8mm and 1mm of verdict, xshift=30mm] (re) {\textbf{RE}};
  \node[b, below right=8mm and 6mm of verdict] (mle) {\textbf{MLE}};
  \node[l, below right=1mm and 2mm of wa.south] (wa1)
    {no parent guard so the rooting double-visits (\probref{c04-reorder-routes}); size
     missing the node itself (\probref{c04-highest-score}); dropped \code{()} flips a lone
     child (\probref{c04-string-tree})};
  \node[l, below right=1mm and 2mm of tle.south] (tle1)
    {a linear \code{find} for the parent each node instead of a \code{kids[]} build
     (\probref{c04-inform}); re-rooting from scratch per node -- take the reroot shift
     (\probref{c04-summit})};
  \node[l, below right=1mm and 2mm of re.south] (re1)
    {recursion on a $2\cdot10^{5}$ chain -- go explicit stack / reverse BFS order
     (\secref{c04-rooting}, Trap~3); no empty-tree guard in a decoder (\probref{c04-recover-preorder})};
  \node[l, below right=1mm and 2mm of mle.south] (mle1)
    {an adjacency \emph{matrix} for a sparse $n{-}1$-edge tree (\secref{c04-reps}); storing
     full token lists when a fused parse-and-place pass fits (\probref{c04-recover-preorder})};
  \draw[flow] (verdict) -- (wa); \draw[flow] (verdict) -- (tle);
  \draw[flow] (verdict) -- (re); \draw[flow] (verdict) -- (mle);
  \draw[flow] (wa) -- (wa1); \draw[flow] (tle) -- (tle1);
  \draw[flow] (re) -- (re1); \draw[flow] (mle) -- (mle1);
\end{tikzpicture}
\caption*{Chapter~4 verdict tree -- only the failures this chapter's problems can
produce.}
\end{figure}

\medskip
\noindent Backup table (same content, terse):
\begin{center}
\small
\renewcommand{\arraystretch}{1.2}
\begin{tabular}{@{}llp{0.42\linewidth}l@{}}
\toprule
\textbf{Verdict} & \textbf{Usual cause} & \textbf{One-line fix} & \textbf{Trained in} \\
\midrule
WA  & bounced to parent      & \code{if (v==p) continue} in both passes & \probref{c04-reorder-routes} \\
WA  & size off by one        & start $\size[u]=1$                      & \probref{c04-highest-score} \\
WA  & lone child on wrong side & keep the empty \code{()}              & \probref{c04-string-tree} \\
TLE & re-root from scratch   & $f(c)=f(p)+(n-2\size[c])$               & \probref{c04-summit} \\
RE  & recursion on a chain   & explicit stack / reverse BFS            & \secref{c04-rooting} \\
MLE & $O(n^2)$ adjacency matrix & a \code{vector<int>} per node        & \secref{c04-reps} \\
\bottomrule
\end{tabular}
\end{center}

\chaptersection{Quick Reference Card}

\begin{quickref}[Chapter 4 at a glance]
\textbf{One question.} A tree is an abstract "who is whose child" relation; before you
compute, you commit it to memory. The commit fixes three things: what a child \emph{is},
whether you can step \emph{up} in $O(1)$, and how many edges you must store
(\defref{c04-forms}).

\medskip
\textbf{Rooting is a view, not data.} Undirected edges carry no above/below; a
parent-guarded DFS manufactures $\parent,\depth,\size$ in one pass, and the result is
\emph{order-independent} (\thmref{c04-dfs-orient}). Moving the root across one edge is an
$O(1)$ shift $f(c)=f(p)+(n-2\size[c])$ (\propref{c04-reroot}) -- that identity is why
``answer for every root'' is $O(n)$, not $O(n^2)$.

\medskip
\textbf{Read the format, then pick the reader.} Level-order array (with nulls) $\to$
\textbf{queue}; preorder with \#/depth marks $\to$ \textbf{stack} or recursion; brackets
$\to$ recursion. A value-only order is uniquely decodable only when nulls are marked
(\thmref{c04-selfdelim}), so when \emph{emitting} a string the empty \code{()} is
load-bearing.
\end{quickref}

\begin{lstlisting}[caption={Chapter 4 code kit: build, root, and read two formats}, label={lst:c04-kit}]
#include <bits/stdc++.h>
using namespace std;
struct TreeNode { int val; TreeNode *left, *right;
    TreeNode(int x=0): val(x), left(nullptr), right(nullptr) {} };

// PARENT ARRAY -> children lists: the one linear idiom.
vector<vector<int>> buildKids(int n, vector<int>& par) {
    vector<vector<int>> kids(n);
    for (int i = 0; i < n; ++i) if (par[i] != -1) kids[par[i]].push_back(i);
    return kids;
}

// UNDIRECTED EDGES -> root: parent/depth (down), size (up). One DFS each pass.
vector<vector<int>> adj; vector<int> par_, dep_, sz_;
void dfsDown(int u, int p) {
    par_[u] = p; dep_[u] = (p == -1 ? 0 : dep_[p] + 1);
    for (int v : adj[u]) if (v != p) dfsDown(v, u);        // THE guard
}
int dfsUp(int u, int p) {
    sz_[u] = 1;                                             // count yourself
    for (int v : adj[u]) if (v != p) sz_[u] += dfsUp(v, u);
    return sz_[u];
}

// REROOT (Summit shape): f(c) = f(p) + (n - 2*size[c]), propagated down.
vector<long long> f_;
void reroot(int u, int p, int n) {
    for (int v : adj[u]) if (v != p)
        { f_[v] = f_[u] + (n - 2LL * sz_[v]); reroot(v, u, n); }
}

// PREORDER WITH DEPTH (dashes) -> tree: a path stack of size == depth.
TreeNode* decodeDashed(const string& s) {                  // LC 1028 style
    stack<TreeNode*> path; TreeNode* root = nullptr; int i = 0, n = s.size();
    while (i < n) {
        int d = 0; while (i < n && s[i] == '-') { ++d; ++i; }
        long long val = 0; while (i < n && isdigit((unsigned char)s[i]))
            { val = val*10 + (s[i]-'0'); ++i; }
        while ((int)path.size() > d) path.pop();
        TreeNode* cur = new TreeNode((int)val);
        if (path.empty()) root = cur;
        else { TreeNode* p = path.top(); if (!p->left) p->left = cur; else p->right = cur; }
        path.push(cur);
    }
    return root;
}
\end{lstlisting}

\chaptersection{Write-From-Memory Drills}

Five signature skeletons, each lifted from a listing already in the chapter. Cover the key
line, reproduce it from memory, then check the answer one line below.

\begin{lstlisting}[caption={Drill 1 -- parent array to children lists (from \coderef{c04-reps})}]
vector<vector<int>> buildKids(int n, vector<int>& par) {
    vector<vector<int>> kids(n);
    for (int i = 0; i < n; ++i)
        if (/* FILL: ______ */) kids[/* FILL: __ */].push_back(i);
    return kids;
}
\end{lstlisting}
\noindent{\scriptsize Key: skip the root (\code{par[i] != -1}); append $i$ to
\code{kids[par[i]]}. One pass, $O(n)$ (\defref{c04-forms}).}

\begin{lstlisting}[caption={Drill 2 -- parent-guarded rooting, down pass (from \coderef{c04-root})}]
void dfsDown(int u, int p) {
    par_[u] = p;
    dep_[u] = (p == -1 ? 0 : /* FILL: ______ */);
    for (int v : adj[u]) if (/* FILL: ______ */) dfsDown(v, u);
}
\end{lstlisting}
\noindent{\scriptsize Key: depth is \code{dep[p]+1}; guard is \code{v != p}. This single
walk is the correct rooting (\thmref{c04-dfs-orient}).}

\begin{lstlisting}[caption={Drill 3 -- bottom-up subtree size (from \coderef{c04-root})}]
int dfsUp(int u, int p) {
    sz_[u] = /* FILL: ___ */;                 // start here!
    for (int v : adj[u]) if (v != p) sz_[u] += /* FILL: ______ */;
    return sz_[u];
}
\end{lstlisting}
\noindent{\scriptsize Key: start at $1$ (a leaf has size $1$); add \code{dfsUp(v,u)}.
Getting this wrong is Trap~2 of \secref{c04-rooting}.}

\begin{lstlisting}[caption={Drill 4 -- reroot shift (from \coderef{c04-summit})}]
void reroot(int u, int p, int n) {
    for (int v : adj[u]) if (v != p) {
        f_[v] = f_[u] + /* FILL: (n - 2*size[v]) */;
        reroot(v, u, n);
    }
}
\end{lstlisting}
\noindent{\scriptsize Key: \code{n - 2*sz[v]} -- the $s$ nodes below get closer, the
$n-s$ above get farther (\propref{c04-reroot}).}

\begin{lstlisting}[caption={Drill 5 -- level-order array to tree (from \coderef{c04-formats})}]
queue<TreeNode*> q; q.push(root);                 // root from vals[0]
int i = 1;
while (!q.empty() && i < n) {
    TreeNode* cur = q.front(); q.pop();
    for (int side = 0; side < 2 && i < n; ++i) {  // left, then right
        if (vals[i] == SENTINEL) continue;        // FILL: a null marks ___ child
        TreeNode* ch = new TreeNode(vals[i]);
        if (side == 0) /* FILL: cur->left  = ch */;
        else           /* FILL: cur->right = ch */;
        q.push(ch);
    }
}
\end{lstlisting}
\noindent{\scriptsize Key: a \code{null}/sentinel slot is an \emph{absent} child, not a
node; the queue makes it breadth-first (\defref{c04-levelorder}, Trap~1).}

\chaptersection{Complete Problem Checklist}

Every problem this chapter offers -- 9 fully solved, 6 practice, one Summit, five mixed.
Tick \emph{only} after an unaided solve; a hinted attempt counts as half. Cite by name and
number -- URLs appear once, in the solved and practice entries.

\begin{center}
\footnotesize
\renewcommand{\arraystretch}{1.25}
\begin{tabular}{@{}c p{0.17\linewidth} l p{0.20\linewidth} p{0.24\linewidth} c@{}}
\toprule
\textbf{\#} & \textbf{Problem} & \textbf{Source} & \textbf{Key technique} & \textbf{What to look for} & \textbf{Done?} \\
\midrule
\multicolumn{6}{l}{\textbf{Section 4.1 -- pointers vs arrays}} \\
W4.1.1 & Max Depth of N-ary Tree & LC 559 & \lv2 children loop & $\max$ child depth $+1$ & $\square$ \\
W4.1.2 & Time to Inform Employees & LC 1376 & \lv3 \code{parent[]}$\to$\code{kids[]} & longest manager chain, memoised & $\square$ \\
W4.1.3 & Validate Binary Tree Nodes & LC 1361 & \lv3 in-degree $+$ reach & one root, every child once, all $n$ & $\square$ \\
P4.1.4 & Collect All Apples in Tree & LC 1443 & \lv3 subtree count & bottom-up, $+2$ per non-empty child & $\square$ \\
P4.1.5 & Time for Tree to Be Infected & LC 2385 & \lv4 build adj $+$ BFS & re-root the fire spread from the start & $\square$ \\
\midrule
\multicolumn{6}{l}{\textbf{Section 4.2 -- rooting an unrooted tree}} \\
W4.2.1 & Reorder Routes to 0 & LC 1466 & \lv3 tagged adjacency & count arrows pointing away from 0 & $\square$ \\
W4.2.2 & Nodes in Sub-Tree Same Label & LC 1519 & \lv4 26-histogram up & merge child counts, self $+1$ & $\square$ \\
W4.2.3 & Minimum Height Trees & LC 310 & \lv4 centroid peel & remove leaves layer by layer & $\square$ \\
W4.2.4 & Nodes With Highest Score & LC 2049 & \lv4 component products & $n{-}\size[u]$ above, $\size[\cdot]$ below & $\square$ \\
P4.2.5 & Most Frequent Subtree Sum & LC 508 & \lv3 post-order sum & hash the running subtree sums & $\square$ \\
P4.2.6 & Longest Path Diff. Chars & LC 2246 & \lv5 top-2 child chains & best two children with $\ne$ label & $\square$ \\
\midrule
\multicolumn{6}{l}{\textbf{Section 4.3 -- encodings and input formats}} \\
W4.3.1 & String from Binary Tree & LC 606 & \lv2 bracket emit & keep \code{()} for a right-only child & $\square$ \\
W4.3.2 & Recover Tree From Preorder & LC 1028 & \lv4 dashed preorder & path stack, pop to depth & $\square$ \\
P4.3.3 & Score of Parentheses & LC 856 & \lv3 depth stack & \code{(} descends, \code{)} folds & $\square$ \\
P4.3.4 & Parse Lisp Expression & LC 736 & \lv4 environment stack & nested scopes on a decode stack & $\square$ \\
\midrule
\multicolumn{6}{l}{\textbf{Summit (chapter end)}} \\
S4 & Sum of Distances in Tree & LC 834 & \lv5 root $+$ reroot & $f(c)=f(p)+(n-2\size[c])$ in $O(n)$ & $\square$ \\
\bottomrule
\end{tabular}
\end{center}

\chaptersection{Recommended Study Order}

Seven focused days. Days 1--2 fix the storage models, days 3--5 the rooting toolkit, days
6--7 the string/format decoders, and day 7 ends on the Summit. Solve in this order even
where it differs from section order -- each day's tool is already in hand from the day
before.

\begin{center}
\small
\renewcommand{\arraystretch}{1.25}
\begin{tabular}{@{}llp{0.55\linewidth}@{}}
\toprule
\textbf{Day} & \textbf{Focus} & \textbf{Tasks} \\
\midrule
1 & Read 4.1 & name the four models on $R$ and $G$; write \code{buildKids} cold; LC 559 \\
2 & Practise 4.1 & LC 1376 (parent$\to$kids), LC 1361 (validity); add LC 1443 as warm-up \\
3 & Read + do 4.2 & root $G$ by hand (\exref{c04-root-run}); LC 1466 tagged edges \\
4 & Practise 4.2 & LC 1519 histogram up; LC 508 subtree sums; LC 2385 build-adj $+$ BFS \\
5 & Centroid \& score & LC 310 leaf-peel; LC 2049 component products \\
6 & Read + do 4.3 & LC 606 emit brackets; LC 1028 dashed-preorder stack; LC 856 depth \\
7 & Summit + revise & LC 834 root-then-reroot; walk the decision flowchart on five inputs \\
\bottomrule
\end{tabular}
\end{center}

\chaptersection{Difficulty Progression}

Tiers are competency gates -- each row names the \emph{skill} the tier unlocks, not just a
rating bucket.

\begin{center}
\small
\renewcommand{\arraystretch}{1.3}
\begin{tabular}{@{}lp{0.36\linewidth}p{0.34\linewidth}@{}}
\toprule
\textbf{Tier} & \textbf{Problems} & \textbf{What passing this tier proves} \\
\midrule
$\lv2$ \ $\bigstar\bigstar$
  & LC 559, 606
  & you can fetch children under any of the four models and emit a bracket string \\
$\lv3$ \ $\bigstar\bigstar\bigstar$
  & LC 1376, 1361, 1443, 1466, 1519, 508, 856
  & a \code{parent[]} or edge list becomes a rooted, aggregated structure reflexively \\
$\lv4$ \ $\bigstar\bigstar\bigstar\bigstar$
  & LC 2385, 310, 2049, 1028, 736, 2458, 2415
  & choose a root to serve the query, decode a marked string, peel a centroid \\
$\lv5$ \ $\bigstar\bigstar\bigstar\bigstar\bigstar$
  & LC 2246, 834 (Summit)
  & reroot an aggregate in $O(n)$; combine a bottom-up return with a top-down pass \\
\bottomrule
\end{tabular}
\end{center}

\chaptersection{Reduction Gallery}

The representation is a disguise; each line collapses a family onto one chapter tool. No
new theory -- just where rooting or decoding is already hiding.

\begin{itemize}[nosep, leftmargin=1.4em]
  \item ``build a tree from \code{parent[]} / \code{leftChild[]},'' (LC 1376, 1361)
        $\longrightarrow$ one loop to \code{kids[]}, then a DFS (\defref{c04-forms}).
  \item ``minimum height / best root'' (LC 310) $\longrightarrow$ the centroid; peel leaves
        from the adjacency (\probref{c04-min-height}).
  \item ``remove a node, score the components'' (LC 2049) $\longrightarrow$ rooted subtree
        sizes: above is $n-\size$, below are child sizes (\probref{c04-highest-score}).
  \item ``reorient edges so all reach 0'' (LC 1466) $\longrightarrow$ root at 0 and tag each
        edge's direction (\probref{c04-reorder-routes}).
  \item ``count a label over each subtree'' (LC 1519, 508, 1443) $\longrightarrow$ one
        bottom-up return of a histogram/sum (\secref{c04-rooting}).
  \item ``sum of distances \emph{from every node}'' (LC 834) $\longrightarrow$ root once,
        then the $O(1)$ reroot shift across each edge (\propref{c04-reroot}).
  \item ``serialise a tree to brackets / a string'' (LC 606) $\longrightarrow$ a preorder
        emit that keeps the null-shape bit (\thmref{c04-selfdelim}).
  \item ``deserialise a dashed / \#-marked preorder'' (LC 1028, 428-style) $\longrightarrow$
        a depth-driven path stack (\probref{c04-recover-preorder}).
  \item ``evaluate a nested bracket language'' (LC 856, 736) $\longrightarrow$ depth is a
        stack; each pair of brackets is a node (\secref{c04-formats}).
  \item ``spread / infection from one node'' (LC 2385) $\longrightarrow$ build the adjacency,
        BFS from the source (\secref{c04-reps}).
\end{itemize}

\chaptersection{Mixed Problem Set}

These refuse to name their pattern; decide the representation, the root, and the reader
before opening the code. A hint tells you the shape of the answer, never the tool.

\begin{mixedset}
  \practice{Height of Binary Tree After Subtree Removal Queries}{LeetCode 2458}
    {https://leetcode.com/problems/height-of-binary-tree-after-subtree-removal-queries/}{LC Medium L4}
    {removing a subtree can only shorten root-to-leaf paths that pass through it -- precompute,
    for every node, the longest and second-longest downward path, then answer a removal from
    the node's depth plus whichever global height survives.}
  \practice{Lowest Common Ancestor of Deepest Leaves}{LeetCode 1123}
    {https://leetcode.com/problems/lowest-common-ancestor-of-deepest-leaves/}{LC Medium L3}
    {you need the deepest leaves' common ancestor without naming depth twice -- let a
    post-order return the height of each side and decide, at every node, whether the deepest
    leaves are on one side or split across both.}
  \practice{Delete Leaves With a Given Value}{LeetCode 1325}
    {https://leetcode.com/problems/delete-leaves-with-a-given-value/}{LC Medium L3}
    {a node only becomes deletable \emph{after} its children have been deleted, so the answer
    must travel back up the call -- what should recursion return so the parent can reattach
    the pruned child?}
  \practice{Reverse Odd Levels of Binary Tree}{LeetCode 2415}
    {https://leetcode.com/problems/reverse-odd-levels-of-binary-tree/}{LC Medium L4}
    {swapping values on every odd depth is a mirror operation -- can you do it with one
    simultaneous walk down the two mirror-image spines, touching a pair once, instead of
    materialising the whole level?}
  \practice{Most Profitable Path in a Tree}{LeetCode 2467}
    {https://leetcode.com/problems/most-profitable-path-in-a-tree/}{LC Medium L3}
    {one player starts at the root and one at a leaf, both racing the same edges -- record,
    for each node, who arrives first and by how much, then score only the nodes the second
    player reaches before the first.}
  \practice{All Paths From Source to Target}{LeetCode 797}
    {https://leetcode.com/problems/all-paths-from-source-to-target/}{LC Medium L3}
    {Read the edge list as an adjacency list first; then every root-to-sink path is
    one DFS with the current walk carried down and released back -- enumerate, don't
    just count.}
  \practice{Is Graph Bipartite?}{LeetCode 785}
    {https://leetcode.com/problems/is-graph-bipartite/}{LC Medium L3}
    {Colour by BFS from every uncoloured start, each neighbour forced to the other
    colour; a clash along one edge names the odd cycle that ruins the split.}
  \practice{Find Eventual Safe States}{LeetCode 802}
    {https://leetcode.com/problems/find-eventual-safe-states/}{LC Medium L4}
    {A node is safe when every walk from it dead-ends at a sink: peel the sinks by
    reverse in-degree, and the peeling order is the topological order over the
    adjacency structure you stored.}
  \practice{Max Network Rank}{LeetCode 1615}
    {https://leetcode.com/problems/max-network-rank/}{LC Medium L3}
    {The rank of a pair is its two degrees minus one if a shared wire exists: the
    answer hides among the top-degree candidates, so only pairs there need an
    adjacency lookup.}
  \practice{Clone Graph}{LeetCode 133}
    {https://leetcode.com/problems/clone-graph/}{LC Medium L4}
    {A deep copy over a cyclic adjacency structure: map every visited node to its
    clone before recursing, so edges rebuilt on the way back already point into the
    copy -- a DFS that carries its visited map as identity.}
\end{mixedset}



\chaptersection{The Summit -- LC 834 Sum of Distances in Tree}

\begin{notebox}[Why this is the chapter's top rope]
Every earlier problem rooted the tree \emph{once} and answered for that one root. The
Summit asks for the answer from \emph{every} node as root -- the naive reading is $n$
traversals, $O(n^2)$. The chapter's single sharpest tool collapses that to $O(n)$: root
once to get subtree sizes, then slide the root across each edge with the $O(1)$ shift of
\propref{c04-reroot}. It is the exact composition this chapter kept promising -- storage
(4.1), one parent-guarded rooting (\thmref{c04-dfs-orient}), and the reroot preview of
\secref{c04-rooting} -- and the reason \probref{c04-summit} was already cited there. The
full rerooting machinery returns in \chref{c20}; here it is earned from scratch.
\end{notebox}

\solvedproblem{c04-summit}{Sum of Distances in Tree}{LeetCode 834}
  {https://leetcode.com/problems/sum-of-distances-in-tree/}{LC Hard L5}
  {root once + subtree sizes, then $O(1)$ reroot shift across every edge (\propref{c04-reroot})}

\begin{problemstatement}[Sum of Distances in Tree]
An undirected tree with $n$ nodes labelled $0,1,\dots,n-1$ is given as $n-1$ edges
\code{edges[i] = [a,b]}. Return an array \code{ans} of length $n$ where \code{ans[i]} is
the sum of the distances from node $i$ to \emph{every} other node (distance is the number
of edges on the unique path). $1 \le n \le 2\cdot10^{4}$.
\end{problemstatement}

\begin{attemptgate}[Attempt it yourself first -- \lv{5}, give yourself 40 minutes]
A fair first attempt: run a BFS/DFS \emph{from each node} and add up distances. That is
correct and $O(n^2)$ -- the whole challenge is noticing that consecutive roots are not
independent.
\peek{if you already know the answer at node $p$ and you move the root one edge to its
child $c$, which nodes get \emph{closer} and which get \emph{farther}?}
\nudge{exactly the two halves split by the edge: the $\size[c]$ nodes in $c$'s subtree drop
by $1$, the other $n-\size[c]$ rise by $1$. Write $f(c)$ in terms of $f(p)$ and
$\size[c]$ -- that is \propref{c04-reroot}.}
\fullhint{Two DFS passes. Pass 1 roots at $0$: compute $\size[u]$ for all $u$ and
$f(0)=\sum_u \depth(u)$. Pass 2 walks down: for each child $c$ of $p$ set
$f(c)=f(p)+(n-2\size[c])$, then recurse. Return $f$. $O(n)$. Full build and proof below.}
\end{attemptgate}

\solstep{Understanding the Problem}
The answer is defined for \emph{each} node as if that node were the root, so a single
rooting is not enough -- but re-rooting from scratch $n$ times is quadratic. The insight is
that the array of answers is not $n$ unrelated numbers: adjacent nodes have adjacent
answers. Pin one root ($0$), harvest the subtree sizes, and then every other answer is a
constant-time edit of its parent's.

\solstep{Key Observations}
\begin{observation}
With the tree rooted at $0$, $f(0)$ is just the sum of all depths -- free from the rooting
pass (\thmref{c04-dfs-orient}), no pairwise distance needed.
\end{observation}
\begin{observation}
Sliding the root across edge $p\!\to\!c$ (where $c$ is $p$'s child under root $0$) moves the
$\size[c]$ nodes of $c$'s subtree one edge \emph{closer} and the $n-\size[c]$ nodes outside
one edge \emph{farther}: $f(c)=f(p)-\size[c]+(n-\size[c])=f(p)+(n-2\size[c])$. This is
\propref{c04-reroot} verbatim.
\end{observation}
\begin{observation}
Because $f$ propagates strictly from parent to child in the rooted tree, one top-down DFS
after pass 1 fills the entire array -- no node is recomputed, so the total is $O(n)$.
\end{observation}

\solstep{Think It Through With Me}
\begin{thinkbox}
\deadend{Run Dijkstra / BFS from every node.}{Correct, but $O(n^2)$ on $n=2\cdot10^4$ is
$4\cdot10^8$ edge relaxations -- a certain TLE. The edges are unweighted and the graph is a
tree, so distances are already given by depth once a root is fixed; there is nothing for a
priority queue to do.}
\deadend{Re-root from scratch: for each candidate root, redo the two-pass DFS.}{$O(n)$ per
root times $n$ roots is again $O(n^2)$. Re-rooting does not \emph{re}compute -- it
\emph{updates} the previous answer across one edge, which is $O(1)$ work.}
\aha{Store $f$ and $\size$. Get $f(0)$ and all $\size$ in one DFS, then one more DFS that
only applies $f(c)=f(p)+(n-2\size[c])$. Two linear passes, every root covered.}
\end{thinkbox}

\solstep{The Approach Ladder}
\begin{center}\renewcommand{\arraystretch}{1.3}
\begin{tabular}{@{}c p{0.235\linewidth} p{0.30\linewidth} l c@{}}
\toprule
\textbf{\#} & \textbf{Approach} & \textbf{Idea} & \textbf{Time} & \textbf{Verdict} \\
\midrule
1 & \lv{2} BFS from every node   & $n$ full traversals          & $O(n^2)$ & \tle \\
2 & \lv{3} DFS from every node   & same count, cheaper constant & $O(n^2)$ & \tle \\
3 & \lv{5} root + reroot shift   & $f(0)$ then $f(c)=f(p)+(n-2\size)$ & $O(n)$ & \ac \\
\bottomrule
\end{tabular}\end{center}
Rungs 1--2 are honest but quadratic -- they re-walk edges the previous root already
covered. Rung 3 spends one extra $\size$ array to convert each new root into $O(1)$ work.

\solstep{Solution}
Root at $0$; in the same DFS compute $\size[u]$ for all $u$ and accumulate $f(0)=\sum
\depth$. Then a second DFS propagates $f$ to every node with \propref{c04-reroot}. The
result array is $f$ (\coderef{c04-summit}).

\solstep{Proof of Correctness}
\begin{proofbox}[Two lemmas, then induction over the rooted tree]
\emph{Lemma 1 (base).} Pass 1 sets $f(0)=\sum_u \depth(u)$. In a tree the distance from the
root $0$ to $u$ \emph{is} $\depth(u)$, so $\sum_u \dist(0,u)=\sum_u\depth(u)=f(0)$ -- the
answer for root $0$ is correct. Pass 1 also computes each $\size[c]$ as the number of nodes
whose path to $0$ passes through $c$ (\thmref{c04-dfs-orient}). \emph{Lemma 2 (step).} For a
parent $p$ and child $c$ under root $0$, moving the root to $c$ changes only the edge
$\{p,c\}$'s contribution: the $\size[c]$ nodes of $\subtree(c)$ lose one edge and the
$n-\size[c]$ nodes outside gain one, so $f(c)=f(p)+(n-2\size[c])$ (\propref{c04-reroot}'s
proof). \emph{Induction.} Pass 2 visits nodes in root-to-leaf order, so whenever it sets
$f(c)$ the parent value $f(p)$ is already correct; by Lemma 2 the child value is then
correct too. The base node $0$ is correct by Lemma 1, hence every $f(v)$ equals
$\sum_u \dist(v,u)$. Each of the $n-1$ edges is applied exactly once (parent to child), so
the array is fully and correctly filled.
\end{proofbox}

\solstep{Implementation}
\begin{lstlisting}[caption={Sum of Distances in Tree (LeetCode 834) -- root once, then reroot}, label={lst:c04-summit}]
#include <bits/stdc++.h>
using namespace std;

class Solution {
    vector<vector<int>> adj;
    vector<int> sz_, ans_;
    int n_;
    // pass 1: root at 0 -> subtree sizes and f(0) (sum of depths)
    void dfs1(int u, int p, int depth) {
        ans_[0] += depth;                 // distance from node 0 is its depth
        sz_[u] = 1;
        for (int v : adj[u]) if (v != p) { // parent guard
            dfs1(v, u, depth + 1);
            sz_[u] += sz_[v];
        }
    }
    // pass 2: propagate f from parent p to child c via the O(1) shift
    void dfs2(int u, int p) {
        for (int v : adj[u]) if (v != p) {
            ans_[v] = ans_[u] + (n_ - 2 * sz_[v]);   // prop:c04-reroot
            dfs2(v, u);
        }
    }
public:
    vector<int> sumOfDistancesInTree(int n, vector<vector<int>>& edges) {
        n_ = n; adj.assign(n, {}); sz_.assign(n, 0); ans_.assign(n, 0);
        for (auto& e : edges) { adj[e[0]].push_back(e[1]); adj[e[1]].push_back(e[0]); }
        dfs1(0, -1, 0);
        dfs2(0, -1);
        return ans_;
    }
};
\end{lstlisting}

\solstep{Trace}
\begin{tracebox}[Running example $G$, rooted at $0$]
Sizes $\size=[7,3,1,2,1,1,1]$, depths $[0,1,1,1,2,2,2]$. Pass 1:
$f(0)=0{+}1{+}1{+}1{+}2{+}2{+}2=9$. Pass 2, top-down:
$f(1)=9+(7-2{\cdot}3)=10$; $f(2)=9+(7-2{\cdot}1)=14$; $f(3)=9+(7-2{\cdot}2)=12$;
$f(4)=10+(7-2{\cdot}1)=15$; $f(5)=10+(7-2{\cdot}1)=15$; $f(6)=12+(7-2{\cdot}1)=17$.
So $\mathit{ans}=[9,10,14,12,15,15,17]$. Check $f(4)$ by hand: distances from $4$ are
$\{2,1,3,3,0,2,4\}$ to $\{0,1,2,3,4,5,6\}$, sum $=15$. \checkmark\ \figref{c04-orient}.
\end{tracebox}

\solstep{Complexity}
\begin{complexitybox}
Time $\bigO{n}$ -- pass 1 and pass 2 each traverse the $n-1$ edges once. Space $\bigO{n}$
for \code{adj}, \code{sz\_}, \code{ans\_} plus $\bigO{h}$ recursion. On the allowed
$n=2\cdot10^{4}$ a chain drives recursion depth to $2\cdot10^{4}$, near the stack limit --
see the hardening box for the iterative rewrite. The answers fit in \code{int} (worst
$\approx n^2/2 = 2\cdot10^{8}$); use \code{long long} if a judge lifts the bound.
\end{complexitybox}

\solstep{Verification and Edge Cases}
\begin{center}\renewcommand{\arraystretch}{1.2}
\begin{tabular}{@{}p{0.30\linewidth} p{0.62\linewidth}@{}}
\toprule
\textbf{Case} & \textbf{What the code does} \\
\midrule
$n=1$ (no edges) & dfs1 adds depth $0$, dfs2 has no child; \code{ans=[0]} \\
Star (one hub) & hub answer $n-1$; each leaf answer $1 + 2(n-2)$ via the shift \\
Path / chain & correct, but recursion is $\bigO{n}$ deep -- go iterative if RE \\
Root label choice & rooting at $0$ is arbitrary; reroot covers all nodes regardless \\
\bottomrule
\end{tabular}\end{center}

\solstep{Pitfalls and Variants}
\begin{warningbox}
Three classic bugs: (i) sizing $\size[c]$ under the \emph{wrong} root -- the sizes must come
from the same root-$0$ orientation used by the shift; (ii) applying $n-2\size[v]$ with
$\size[v]$ read for the wrong endpoint (it is the \emph{child's} subtree); (iii) a plain
\code{int} overflow if the bound is lifted past $2\cdot10^{4}$. Variant -- the same two-pass
shape answers ``count nodes where $\size[u]$ divides $n$'', ``product of component sizes on
removal'' (LC 2049, \probref{c04-highest-score}), and every reroot DP in \chref{c20}: only
the definition of $f$ and the shift change.
\end{warningbox}

\begin{insightbox}[Key takeaways -- three things to remember]
\begin{enumerate}[nosep, leftmargin=1.6em]
  \item ``Answer for every root'' is \emph{not} $n$ traversals -- a root change across one
        edge is an $O(1)$ edit (\propref{c04-reroot}).
  \item The whole trick needs one number per node, $\size[c]$, which the \emph{same} rooting
        DFS that gave $f(0)$ already computed (\thmref{c04-dfs-orient}).
  \item Rooting is a \emph{view}: pick root $0$ for convenience, then let the shift move the
        viewpoint. This chapter's three topics compose into a single $O(n)$ machine.
\end{enumerate}
\end{insightbox}

\begin{hardeningbox}
\begin{itemize}[nosep, leftmargin=1.4em]
  \item \textbf{RE on a chain}: replace both DFS with an explicit stack -- build a parent
        order by an iterative DFS, process it in \emph{reverse} to accumulate $\size$, then
        walk it \emph{forward} to apply the reroot shift. Same $O(n)$, zero call frames.
  \item \textbf{Answer type}: compute $f$ and $\size$ as \code{long long} if $n$ can exceed
        $2\cdot10^{4}$; the intermediate $n - 2\size[v]$ can also go negative for a large
        child, so keep it signed.
  \item \textbf{Generalise}: swap the shift for any per-edge delta to reroot a different
        aggregate (paths, maxima, colour counts). The full catalog of which aggregates admit
        an $O(1)$ shift is the Rerooting chapter (\chref{c20}).
\end{itemize}
\end{hardeningbox}

\chaptersection{Tool Kit Inventory}

The machine assembled so far -- cumulative through Chapter~4. Chapters~1--3's rows are
copied forward unchanged; Chapter~4's rows are appended. Pointers only, never prose.

\begin{center}
\small
\renewcommand{\arraystretch}{1.25}
\begin{tabular}{@{}p{0.24\linewidth} l c p{0.27\linewidth} p{0.16\linewidth}@{}}
\toprule
\textbf{Tool} & \textbf{Since} & \textbf{Level} & \textbf{Trigger phrase} & \textbf{Where} \\
\midrule
\multicolumn{5}{l}{\itshape Chapters 1--3 -- carried forward} \\
Adjacency list + parent-guard DFS & 1.1 & \lv1 & edge list given, not pointers & \coderef{c01-build} \\
Top-down carry / bottom-up return & 1.2 & \lv1 & value from root / subtree aggregate & \coderef{c01-measure} \\
Structural skeleton (null base + 2 calls) & 2.1 & \lv1 & any node-plus-two-children question & \coderef{c02-skeleton} \\
Return-new-root transform & 2.4 & \lv3 & the pointer graph itself must change & \coderef{c02-transform-skeleton} \\
Three-slot DFS (pre/in/post) & 3.1 & \lv1 & a node sequence in a named order & \coderef{c03-dfs} \\
Explicit-stack traversal / Morris & 3.2 & \lv2--\lv4 & recursion depth under attack & \coderef{c03-inorder-stack} \\
Level-order BFS with frontier snapshot & 3.3 & \lv2 & output or reduce \emph{per depth} & \coderef{c03-bfs} \\
Pair reconstruction / marked codec & 3.4 & \lv4--\lv5 & two traversals given / round-trip & \coderef{c03-build} \\
\midrule
\multicolumn{5}{l}{\itshape Chapter 4 -- new this chapter} \\
Four storage models + children fetch & 4.1 & \lv1 & commit edges to memory; how to reach kids & \coderef{c04-reps} \\
Parent array $\to$ children lists & 4.1 & \lv2 & a \code{parent} array, go down & \coderef{c04-reps} \\
Root: read parent, depth, size, height & 4.2 & \lv3 & undirected edges, need above/below & \coderef{c04-root} \\
Centroid leaf-peel (best root) & 4.2 & \lv4 & minimise height / choose a root & \coderef{c04-min-height} \\
Level-order decode (queue) & 4.3 & \lv2 & input is LeetCode's null-filled array & \coderef{c04-formats} \\
Bracket / serialise emit & 4.3 & \lv2 & tree $\to$ self-delimiting string & \probref{c04-string-tree} \\
Depth-preorder decode (path stack) & 4.3 & \lv4 & dashed or \#-marked preorder & \coderef{c04-recover-preorder} \\
Reroot shift $f(c)=f(p)+(n-2\size)$ & Summit & \lv5 & the same aggregate from \emph{every} root & \coderef{c04-summit} \\
\bottomrule
\end{tabular}
\end{center}

\chaptersection{Related Topics}

\begin{itemize}[nosep, leftmargin=1.4em]
  \item \chref{c01}, \chref{c02} -- backward: the adjacency-build loop \coderef{c01-build}
        and the \code{TreeNode} model this chapter now treats as one of four interchangeable
        storages.
  \item \chref{c03} -- backward: every rooting/decoding pass here \emph{uses} the DFS orders;
        this chapter decides where the children live, not when you act on them.
  \item \chref{c24} -- Euler tour / \code{tin},\code{tout}: the next step -- turn a subtree
        into a contiguous array range so it can carry a Fenwick/segment tree.
  \item \chref{c20} -- Rerooting: \probref{c04-summit}'s $O(1)$ shift, generalised to any
        aggregate admitting a parent$\to$child delta.
  \item \chref{c11} -- LCA and lifting: root once here, then binary lifting answers
        $\dist(u,v)$ and ancestor queries for \emph{pairs} of nodes.
  \item \chref{c07}, \chref{c05}, \chref{c32} -- the heap's implicit array is
        \thmref{c04-implicit} specialised to a complete tree; BSTs fix an order on pointers;
        a Prufer sequence is a sixth encoding (\secref{c04-formats}).
\end{itemize}

\begin{lastword}[Last Word -- Chapter 4]
\emph{``A tree is an idea; a representation is a decision. Every algorithm is only as fast
as the navigation you paid for.''}

\medskip
Chapters~2 and~3 asked \emph{what} a node hands back and \emph{when} you act. This chapter
asked the more physical question: \emph{where do the edges live}. The answer -- pointer,
adjacency, parent array, implicit index -- is a trade you make once and then live with,
because it fixes whether ``go to my parent'' is $O(1)$ or $O(n)$ and how many of the
$n-1$ edges you must store (\defref{c04-forms}). Rooting then revealed that ``above/below''
is not a fact about the tree but a \emph{view} you install in one DFS and can move across an
edge in constant time (\propref{c04-reroot}).

\medskip
Carry three reflexes forward:
\begin{enumerate}[nosep, leftmargin=1.6em]
  \item \emph{Read the input format before writing anything.} An edge list, a
        \code{parent[]}, a dashed string, a null-filled array -- each names its own reader
        (adjacency $+$ guard, one build loop, a stack, a queue).
  \item \emph{Root once, then move the root.} A parent-guarded DFS gives $\parent/\depth/\size$
        for free; the reroot shift turns ``every root'' from $O(n^2)$ into $O(n)$.
  \item \emph{Guard the shape bit.} Recursion on a chain overflows, and a serial format
        without a null marker is many-to-one -- keep the \code{()} and keep a stack
        (\thmref{c04-selfdelim}).
\end{enumerate}

\medskip
And three habits, drilled until boring: write $G$'s $\parent$, $\size$ and bracket encoding
by hand before trusting a solution; keep one sentence straight -- \emph{the parent array
answers ``up'' for free and ``down'' after one build loop}; and never decode a level-order
array with a preorder recursion, because the array was written breadth-first.

\medskip
\textbf{Next:} \chref{c24} keeps the array you now know how to build and asks a sharper
question -- lay the tree out along a line so that \emph{``the whole subtree of $v$''}
becomes one contiguous interval $[\code{tin}[v],\code{tout}[v]]$, and a subtree query
becomes a range query you already know how to answer.

\begin{center}\textcolor{softpurple}{\rule{0.4\linewidth}{0.6pt}\quad $\blacksquare$\quad \rule{0.4\linewidth}{0.6pt}}\end{center}
\end{lastword}
